\documentclass[10pt]{article}
\usepackage[utf8]{inputenc}
\usepackage[english]{babel}
\usepackage{amsmath}
\usepackage{cite}
\title{The Riemann Hypothesis: Distribution of Prime Numbers\\and the Zeros of the Zeta Function}

\date{\today}
\begin{document}
\maketitle

\paragraph{}
The distribution of prime numbers is one of the oldest and deepest 
problems in mathematics. Gauss and Legendre conjectured around 1800 
that the number of primes up to $x$, denoted $\pi(x)$, behaves 
asymptotically as $x / \ln x$. This result, known as the Prime Number 
Theorem, was proved in 1896 by Hadamard and de la Vallée Poussin.
In 1859, Riemann introduced the zeta function, defined for $\Re(s) > 1$ by
\[
    \zeta(s) = \sum_{n=1}^{\infty} \frac{1}{n^s} = \prod_{p \text{ prime}} \frac{1}{1-p^{-s}}
\]
and extended it to the whole complex plane by analytic continuation.
Riemann observed that $\zeta(s)$ vanishes at all negative even integers 
(trivial zeros), and that all remaining non-trivial zeros lie in the 
strip $0 < \Re(s) < 1$. He conjectured that all non-trivial zeros lie 
exactly on the critical line $\Re(s) = 1/2$. This conjecture, known as 
the Riemann Hypothesis, remains unproved after more than 160 years and 
is listed as one of the seven Millennium Prize Problems 
\cite{carlson_millennium_2023}. A general overview of the properties 
of $\zeta(s)$ is given in \cite{gourdon_riemann_nodate}.

\paragraph{}
Hardy proved in 1914 that infinitely many zeros lie on the critical line. 
This was later strengthened by Hardy and Littlewood (1921), who obtained 
quantitative estimates. Selberg (1942) showed that a positive proportion 
of zeros lie on the critical line, and Conrey pushed this proportion above 
two fifths \cite{noauthor_more_1989}, which remains the best known result 
to this day. On the computational side, Odlyzko and Schonhage developed 
fast algorithms based on the Riemann-Siegel formula to evaluate $\zeta(s)$ 
efficiently at many points, enabling large-scale numerical verification 
\cite{odlyzko_fast_nodate}. These computations have confirmed that the 
first $10^{22}$ zeros all lie on the critical line, providing strong 
empirical support for the hypothesis. Another major direction connects 
the zeros of $\zeta(s)$ to quantum mechanics and random matrix theory: 
the statistical spacing between consecutive zeros follows the same 
distribution as eigenvalues of random unitary matrices. This unexpected 
connection, supported by the numerical work of Odlyzko, suggests that 
the zeros might be eigenvalues of some unknown Hermitian operator 
\cite{noauthor_riemann_nodate}. This is known as the Hilbert-Pólya 
conjecture and has inspired a rich interdisciplinary research program 
at the intersection of number theory and mathematical physics.

\paragraph{}
Despite these advances, the Riemann Hypothesis remains one of the most 
important open problems in mathematics \cite{carlson_millennium_2023}. 
Several directions are actively pursued. First, the proportion of zeros 
proven to lie on the critical line is still far from 100\%, and pushing 
this bound further requires new analytic tools such as improved mollifiers 
and mean-value estimates for $\zeta(s)$. Second, current zero-free regions 
are far from the critical line, and stronger results would yield better 
error terms in the Prime Number Theorem unconditionally. Third, the 
Hilbert-Pólya conjecture remains open: no Hermitian operator whose 
eigenvalues are the imaginary parts of the zeros has been constructed 
rigorously \cite{noauthor_riemann_nodate}. A deeper understanding of 
the analytic properties of $\zeta(s)$ \cite{gourdon_riemann_nodate} 
and more powerful tools from analytic number theory 
\cite{iwaniec_analytic_2021} are needed to make further progress 
toward a proof.

\bibliographystyle{plain}
\bibliography{references}
\end{document}